% BÀI TOÁN HÌNH HỌC - BAI_HINH_31 (TRÍCH ĐỀ 31)
\begin{ex}
Cho tam giác ABC nhọn $AB < AC$ có đường cao AD và đường phân giác trong AO (D, O thuộc cạnh BC). Kẻ OM vuông góc với AB tại M, ON vuông góc với AC tại N.
	\begin{enumerate}
		\item[a)] Chứng minh bốn điểm A, D, O, N cùng nằm trên một đường tròn.
		\item[b)] Chứng minh: $\widehat{BDM} = \widehat{ODN}$. Qua O, kẻ đường thẳng vuông góc với BC cắt MN tại I, AI cắt BC tại K. Chứng minh K là trung điểm của BC.
	\end{enumerate}
\end{ex}
